\section{Цель работы}
\textbf{Цель работы:} экспериментальное исследование и анализ погрешностей позиционирования навигационной аппаратуры потребителей (НАП) глобальных навигационных спутниковых систем (ГНСС) в различных городских условиях.

\section{Постановка задачи}
\textbf{Для достижения поставленной цели необходимо решить следующие задачи:}
\begin{enumerate}
    \item Провести эксперимент по определению местоположения с использованием сервиса геолокации смартфона (в качестве навигационного сервиса выбрано приложение Яндекс Карты, демонстрирующее стабильную работу в районе проведения исследования и GPtest для определения колличества спутников) в трех различных условиях:
    \begin{itemize}
        \item на открытой местности (в радиусе 50 метров от которой нет зданий);
        \item вблизи малоэтажной застройки (здания от 2 до 4 этажей);
        \item вблизи многоэтажной застройки (здания более 9 этажей).
    \end{itemize}
    \item Обеспечить корректный захват сигналов, выжидая 3--5 минут в каждой локации.
    \item Вычислить ошибку позиционирования для каждого сценария путем сравнения истинного положения с координатами, определенными приемником, используя встроенный инструмент «Линейка».
    \item Задокументировать результаты измерений с помощью снимков экрана приложения и фотографий (селфи) на фоне исследуемых локаций.
    \item Провести сравнительный анализ полученных данных о погрешностях и сделать выводы о влиянии высотности и плотности застройки на точность позиционирования НАП ГНСС.
\end{enumerate}

Все измерения производились в дневное время суток, в переменную облачность, что может сказаться на точности определения местоположения.